\section{职业与社会：Builder 身份、短期阵痛与长期收益}

关于“程序员是否被替代”，Peter 的答案并不温和：纯手工编码的稀缺性会下降。但他强调，构建能力并不会消失，只会重新分层。

\subsection{“编程像 knitting”这句话该怎么读}

这句比喻容易引发反感，但其技术含义是：\textbf{实现层的边际价值下降，问题定义与系统设计层的价值上升}。不是“不需要人”，而是“需要的人做不同的事”。

\begin{knowledgebox}{Builder 身份的迁移路径}
\begin{itemize}
    \item 从“语法掌握”迁移到“问题建模”；
    \item 从“代码产量”迁移到“系统结果负责”；
    \item 从“单人实现”迁移到“人机协作编排”。
\end{itemize}
\end{knowledgebox}

\subsection{社会层面的双重现实}

访谈末段同时给出两类信号：一是就业和身份焦虑真实存在；二是小企业自动化与残障用户赋能也真实发生。两者必须同时成立，讨论才完整。

\begin{importantbox}{政策与产品都绕不开的现实}
\begin{enumerate}
    \item 短期结构性失业和技能错配会加剧；
    \item 长期生产力提升会创造新岗位，但时间滞后；
    \item 教育体系需从“工具教学”转向“系统思维教学”；
    \item 企业责任从“提高效率”扩展到“降低转型痛苦”。
\end{enumerate}
\end{importantbox}

\begin{figure}[H]
\centering
\includegraphics[width=0.9\textwidth]{frames/frame-025932.jpg}
\caption{平台冲突与用户诉求：Agent 访问权成为新博弈点\protect\footnotemark}
\end{figure}
\footnotetext{视频画面时间区间：02:58:20--03:00:10。}

\subsection{本章小结}
OpenClaw 的社会争议提醒我们：技术叙事若只谈“会更好”，会失去可信度。严肃讨论必须同时面对效率增益与转型代价。

